\documentclass[10pt,a4paper]{amsart}
\usepackage[margin=19mm]{geometry}
\usepackage{amsmath,amssymb,mathtools}
\usepackage[scr=rsfs,cal=euler]{mathalfa}
\usepackage{xfrac}
\usepackage[unicode,hidelinks,bookmarks=false]{hyperref}
\newcommand{\ie}{i.e.\ }
\newcommand{\cf}{cf.\ }
\newcommand{\resp}{respectively\ }
\newcommand{\Subsection}{\subsection}
\providecommand{\nobreakdash}{\nobreak\hbox{-}}
\setlength{\emergencystretch}{2em}
\multlinegap0pt
\usepackage{amssymb}
\usepackage{enumerate}
\usepackage{csquotes}
\usepackage{mathtools}
\usepackage{tikz}
\newtheorem{theorem}{Theorem}
\title{Some inequalities related to Heinz mean constant with Birkhoff orthogonality}
\author{Kallal Pal, Sumit Chandok}
\date{Source: arXiv:2512.24675v1 (CC BY 4.0)}
\begin{document}
\maketitle

\noindent\emph{Source and licence.} Some inequalities related to Heinz mean constant with Birkhoff orthogonality. Version arXiv:2512.24675v1 (CC BY 4.0), distributed by arXiv under the Creative Commons Attribution 4.0 International licence, http://creativecommons.org/licenses/by/4.0/, as recorded in the arXiv licence metadata for this exact version and shown on the abstract page https://arxiv.org/abs/2512.24675v1. Full licence notice: \texttt{licenses/arxiv-2512.24675-CC-BY-4.0.txt}.\par\smallskip

\noindent\textit{Editorial scope.} Counted unit: the article's theorem t11 (source0.tex lines 430--432). The article's complete original proof of it is delivered with it (source0.tex lines 433--515, one \texttt{proof} environment of 83 lines). No mathematics is written, repaired or re-derived: the statement and the proof are the article's own lines, in the article's own notation, and no retained line is substituted or re-flowed.  No further source-local context is needed: every cross-reference of every retained line resolves inside the two delivered spans.  Nothing was omitted from any retained span, so the package carries no omission record.  No retained line contains a citation, so the wrapper carries no bibliography and no bibliography entry.  No retained line contains a figure environment, an image-inclusion command or drawing code, so no image is delivered and the package declares no asset dependency.  Editorial formatting only, in addition: an AMS \texttt{amsart} preamble that restates the article's own declarations and macros used by the retained lines, namely the environments \texttt{theorem} on separate counters, together with the standard packages those lines need.  The retained environments print in this document as Theorem 1 in order of appearance, whereas the article prints the same items with the numbering of its own full document.  The complete native source of the article as distributed in the arXiv e-print for this exact version is bundled unchanged as \texttt{sources/191918/source0.tex}.
\begin{theorem}\label{t11}
A Banach space $X$ with $H_\nu(X,B)< 2$,  is uniformly non-square.
\end{theorem}

\begin{proof}
On the contrary, we suppose that $X$ is not uniformly non-square.
Then there exist $x,y \in S_X$ such that
\[
\|x \pm y\| \ge 2-\varepsilon^2, ~\text{for}~ \varepsilon \in (0,1/2).
\]
On rewriting we obtain
\[
2-\varepsilon^2 \le \|x+y\|
\le \|x+\varepsilon y\| + \|(1-\varepsilon)y\|
= \|x+\varepsilon y\| + 1-\varepsilon.
\]
Hence
$\|x+\varepsilon y\| \ge 1+\varepsilon-\varepsilon^2 > 1$. Also, $\|x+\varepsilon y\|\leq 1+\varepsilon$.

Define $\varpi_2=\frac{\|x+\varepsilon y\|-1}{\varepsilon}x-y$, 
$\bar \varpi_2=\frac{\varpi_2}{\|\varpi_2\|}$,
and $\phi:[0,1]\to\mathbb{R}$ by $\phi(\alpha)=\|y+\alpha \varpi_2\|$.
It is easy to see that $\phi$ is convex and 
$\phi\!\left(\frac{1}{\|x+\varepsilon y\|}\right)
=\frac{\|x+\varepsilon y\|-1}{\varepsilon}
=\phi(1)$.
Thus, $\phi$ attains its minimum at some
$t\in\left[\frac{1}{\|x+\varepsilon y\|},\,1\right]
\subset\left[{1+\varepsilon},1\right]
\quad (\text{since } \|x+\varepsilon y\|\le 1+\varepsilon)$,
which implies that $y+t \varpi_2 \perp_B \varpi_2$.

Choose $\varpi_1=y+t \varpi_2$ and $\bar \varpi_1=\varpi_1/\|\varpi_1\|$.
Since Birkhoff orthogonality is homogeneous, $\bar \varpi_1\perp_B \bar \varpi_2$ with $\|\bar \varpi_1\|=\bar \varpi_2\|=1$.

By the convexity of the norm,
$1-\varepsilon^2 \le \|s x \pm (1-s)y\| \le 1$, for all  $s\in[0,1]$. It implies that
$\|\alpha x \pm \beta y\|
\in \bigl[(1-\varepsilon^2)(\alpha+\beta),\,\alpha+\beta\bigr]$
for all non-negative scalars $\alpha,\beta$.\\
% Notice that
% \[
% 1-\varepsilon
% =\frac{1+\varepsilon-\varepsilon^2-1}{\varepsilon}
% \le \frac{\|x+\varepsilon y\|-1}{\varepsilon}
% \le \frac{1+\varepsilon-1}{\varepsilon}
% =1.
% \]
Hence,
\begin{align*}
\|\varpi_2\|&
=\left\|\frac{\|x+\varepsilon y\|-1}{\varepsilon}x-y\right\|
\in \bigl[(1-\varepsilon^2)(2-\varepsilon),\,2\bigr],\\
\|\varpi_1\|&
=\left\|t\frac{\|x+\varepsilon y\|-1}{\varepsilon}x+(1-t)y\right\|
\in \bigl[(1-\varepsilon^2)(1-\varepsilon),\,1\bigr],\\
\left\|\varpi_1+\frac{\varpi_2}{2}\right\|&
=\left\|\left(t+\frac12\right)\frac{\|x+\varepsilon y\|-1}{\varepsilon}x
-\left(t-\frac12\right)y\right\|
\ge (1-\varepsilon^2)\left(\frac{2}{1+\varepsilon}-\frac{3\varepsilon}{2}\right),\\
&\hspace{-1cm}\text{and}~~
\left\|\varpi_1-\frac{\varpi_2}{2}\right\|\geq \|u\|\geq (1-\epsilon^2)(1-\epsilon).
\end{align*}
Therefore
\begin{align*}
\|\bar \varpi_1+\bar \varpi_2\|
&\ge \left\|\varpi_1+\frac{\varpi_2}{2}\right\|
 - \|\varpi_1-\bar \varpi_1\| - \left\|\bar \varpi_2-\frac{\varpi_1}{2}\right\| \\
&= \left\|\varpi_1+\frac{\varpi_2}{2}\right\| - (1-\|\varpi_1\|) - \frac{2-\|\varpi_2\|}{2} \\
&\ge \frac{2-2\varepsilon^2}{1+\varepsilon} - 3\varepsilon-2\epsilon^2\geq \frac{2-2\varepsilon^2}{1+\varepsilon} - 5\varepsilon\\
&=\frac{2-5\varepsilon-7\varepsilon^2}{1+\varepsilon}.
\end{align*}
and 
\begin{align*}
 \|\bar \varpi_1-\bar \varpi_2\| &\ge \left\|\varpi_1-\frac{\varpi_2}{2}\right\|
 - \|\varpi_1-\bar \varpi_1\| - \left\|\bar \varpi_2-\frac{\varpi_2}{2}\right\| \\  
&= \left\|\varpi_1-\frac{\varpi_2}{2}\right\| - (1-\|\varpi_1\|) - \frac{2-\|\varpi_2\|}{2} \\ 
&\geq 1-\frac{5\varepsilon}{2}-3\varepsilon^2+\varepsilon^3\\
&\geq 1-4\varepsilon.
\end{align*}
Using the definition, we get
\begin{align*}
H_\nu(X,B)&\geq\frac{\|\bar \varpi_1+\bar \varpi_2\|^{\nu}\|\bar \varpi_1-\bar \varpi_2\|^{1-\nu} + \|\bar \varpi_1+\bar \varpi_2\|^{1-\nu}\|\bar \varpi_1-\bar \varpi_2\|^{\nu}}{2}\\
&\geq \frac{(\frac{2-5\varepsilon-7\varepsilon^2}{1+\varepsilon})^\nu(1-4\epsilon)^{1-\nu}+(\frac{2-5\varepsilon-7\varepsilon^2}{1+\varepsilon})^{1-\nu}(1-4\epsilon)^\nu}{2}.
\end{align*}
Letting $\varepsilon\rightarrow 0$, we get $H_\nu(X,B)\geq 2$, a contradiction.
\end{proof}

\end{document}
